\section{Detailed Goal-Proxy Assumptions}

The benchmark treats a disappearing goal as a task-identity problem, not only as a missing-pixels problem. A referent can disappear because the robot arm occludes it, because the manipulated object blocks the view, because the goal is inside a container, because a human moves between the camera and the goal, or because the robot must execute after the original observation is stale. In each case, the robot must decide whether the old goal identity is still valid.

A goal proxy is a compact planning object. It stores the identity binding, the geometric anchor, the task role, the affordance, an uncertainty state, a re-binding gate, and a re-acquisition policy. The identity binding says which physical or semantic referent the task is about. The geometric anchor says where action should be applied. The task role says why this object or region is the goal. The re-binding gate says which new observations are compatible with the old goal.

This separation matters because a robot can have good geometry and bad identity. A manipulator may know that there is an object at a location and still act on the wrong object. It can also have weak geometry and good identity. A robot may know the intended cup but need active perception before final placement. The final paper therefore evaluates success, final error, identity swaps, proxy loss, re-binding precision, re-acquisition success, search effort, delay, waste, safety exposure, progress, and utility separately.

\section{Proxy State And Update Contract}

A deployed proxy interface should expose a small state tuple:
\[
z_g = (i_g, r_g, x_g, a_g, u_g, b_g, \tau_g),
\]
where $i_g$ is identity, $r_g$ is task role, $x_g$ is geometry, $a_g$ is affordance, $u_g$ is uncertainty, $b_g$ is the re-binding gate, and $\tau_g$ is a stale-age or disappearance timer. The action policy consumes $z_g$ rather than a raw goal observation.

The update contract has three rules. First, visibility updates geometry but does not automatically change identity. Second, disappearance increases uncertainty and stale age but does not erase the proxy. Third, reappearance must pass a gate before it can update identity. If the gate fails, the policy either searches, requests confirmation, escalates, or aborts.

Loose proxy acceptance violates the third rule. It lets a nearby or appearance-similar object update the proxy without enough evidence. The v2 stress and the full-scale ambiguity table both show why that is dangerous. A proxy without a gate can be worse than a last-seen belief because it gives the planner a false sense of continuity.

\section{Risk-Calibrated Identity Proxy}

\method estimates identity risk from ambiguity, false-observation risk, hidden motion, human-proximity risk, action irreversibility, diagnosis quality, and semantic-anchor strength. The benchmark uses a deterministic score:
\[
\rho = \alpha a + \beta f + \gamma d + \delta h + \kappa c - \eta q - \mu s,
\]
where $a$ is ambiguity, $f$ is false-observation risk, $d$ is drift or hidden motion, $h$ is human risk, $c$ is action cost, $q$ is diagnosis quality, and $s$ is semantic-anchor strength.

The score controls five decisions:
\begin{enumerate}[leftmargin=*]
\item continue acting on the proxy;
\item tighten the re-binding gate;
\item search actively with wrist or viewpoint motion;
\item request tactile, geometric, or semantic confirmation;
\item stop or escalate if the identity risk is too high.
\end{enumerate}

The method is intentionally simple. The contribution is not a learned re-identification network. It is a policy contract: a manipulation planner should report identity risk and spend re-acquisition effort only when the evidence supports it. A stronger learned version can replace the score, but it should still be evaluated with the same metrics.

\section{Operational Load}

\method is not the most aggressive search policy. Active-search proxy has aggregate search effort 0.756 and utility 0.453. \method has search effort 0.309 and utility 0.595. The gap shows why search should be calibrated. Searching always is safer on identity swaps, but it wastes time and delays action in easy regimes.

Proxy with re-acquisition is also strong. It has success 0.762, swap 0.145, re-acquisition success 0.586, and utility 0.454. It often finishes tasks more often than \method, but it accepts more identity risk and spends more search effort. The proposed method wins because it balances task progress, precision, swap risk, and search cost.

Visible-only has low swap rate, 0.020, but low success, 0.284, and high proxy loss, 0.496. This baseline is useful because it exposes the other extreme: avoiding identity swaps by refusing to act after disappearance is not enough for manipulation.

\section{Ambiguity Findings}

The ambiguity stress is the scientific center of the paper. Loose proxy acceptance has swap rate 0.249 even under clear reappearance. It rises to 0.466 with close spatial distractors, 0.509 with appearance-similar distractors, 0.592 with crossing trajectories, 0.605 when the goal moves while hidden, and 0.705 under adversarial false observations.

Proxy re-acquisition reduces swaps, but it is not perfect. The aggregate swap rate under adversarial false observation remains 0.190. \method further reduces that to 0.106 while keeping utility 0.542. That is the final claim in compact form: persistent goal proxies need identity-calibrated re-binding, not silent acceptance.

The table also shows graceful degradation. \method utility falls from 0.662 under clear reappearance to 0.542 under adversarial false observation. That decline is important. A paper that reports a flat result under ambiguity would be suspicious. The method should degrade because the task is genuinely harder.

\section{Horizon Findings}

Disappearance horizon exposes the weakness of visible-only and last-seen baselines. Visible-only success falls from 0.358 under short disappearance to 0.223 under mixed long horizons. Last-seen error rises from 0.219 to 0.349, and its swap rate rises from 0.345 to 0.431.

\method remains useful across horizons, but it also degrades. Its utility falls from 0.868 under short disappearance to 0.801 under mixed long horizons in the original high-success scoring run, and to the lower calibrated values in the final utility table after scoring was tightened. The final manuscript reports the generated final table values. The qualitative pattern is the point: long disappearance requires identity maintenance and re-acquisition, not only memory.

The horizon axis is especially relevant for real manipulation. The decisive action often happens after the target leaves view. A shelf slot may be visible during approach but hidden at insertion. A cup may be visible before the wrist occludes it. A drawer interior may disappear as the robot reaches in. The planner must preserve the goal through that interval.

\section{Occlusion Findings}

Different occlusion causes demand different re-acquisition strategies. Robot self-occlusion can often be solved by a wrist camera or brief viewpoint change. Container-interior disappearance may require geometric reasoning and contact constraints. Human occlusion requires conservative action because the same disappearance can imply safety risk.

The benchmark therefore separates occlusion from ambiguity. A goal can be heavily occluded but unambiguous, or lightly occluded but ambiguous because a distractor is nearby. Reporting only occlusion length would hide this distinction. The final paper reports both axes.

Visible-only proxy loss is high across occlusion regimes. Re-acquisition spends more search under hard occlusions. \method keeps swap risk bounded by increasing caution when occlusion and ambiguity combine.

\section{Observability And Re-Acquisition}

Passive vision-only reappearance is the weakest observability regime. It gives the robot whatever happens to become visible and asks the proxy gate to decide whether to accept it. Active wrist-camera motion improves diagnosis but costs time. Tactile/geometric confirmation gives strong identity evidence but may require contact. Semantic anchors help language-specified goals but can be weak when several objects share the same label. Multimodal delayed confirmation is strongest but slowest.

The benchmark uses these regimes to show that re-acquisition is not one thing. It can be visual, tactile, geometric, semantic, or multimodal. A submission-ready disappearing-goal paper should report which kind of evidence re-bound the proxy.

\method uses re-acquisition when identity risk is high and search cost is justified. Active-search proxy uses it more aggressively. That is why active search has high re-acquisition success but loses utility in costly regimes.

\section{Task-Level Findings}

The task table shows that the method is not tuned to one toy scenario. Utility remains positive across tabletop cup placement, shelf insertion, drawer placement, human handover, bimanual assembly, hidden fasteners, cable routing, bin identity, mobile view loss, language referents, regrasping, and failed-grasp recovery.

The hardest tasks are cable aperture routing, mobile manipulation view loss, and human handover. Cable and mobile manipulation combine geometry drift with disappearance. Handover adds safety risk. The easiest tasks are tabletop cup placement and language referent clearing because their geometry or semantic anchors are less adversarial in the generator.

These task findings are not real-robot claims. They are a coverage audit. They show that the proposed mechanism was not evaluated only on the easiest disappearing-goal case.

\section{V2 Negative Control}

The v2 stress remains in the final paper because it prevents overclaiming. In close-distractor reappearance, loose proxy acceptance had success 0.094 and swap rate 0.906. In severe ambiguity, loose proxy acceptance had success 0.020 and swap rate 0.980. Re-acquisition improved those numbers to 0.928 and 0.821 success, but with nonzero swap rates.

This result changes the final thesis. The paper does not claim that persistent proxies solve disappearing goals. It claims that persistent proxies are useful only when identity re-binding and re-acquisition are explicit and audited. The full-scale benchmark scales that lesson across many tasks and regimes.

\section{Metric Semantics}

Mission success measures whether the task completes without violating the goal identity. Final error measures geometric distance from the true goal. Identity swap measures whether the robot acted on the wrong referent. Proxy loss measures whether the policy lost a usable goal object. Re-binding precision measures whether accepted observations were compatible with the original goal. Re-acquisition success measures whether active information gathering recovered the goal.

Search effort measures active perception, tactile checks, or semantic confirmation. Delay measures time-like cost from waiting or searching. Waste measures unnecessary search or false confirmation work. Safety exposure measures human-proximity or irreversible-action risk. Progress measures how well the robot continues moving toward the intended task during disappearance.

Utility combines these components. It penalizes swaps strongly because an identity swap is not just geometric error; it is a task change. It also penalizes proxy loss, search, delay, waste, and safety exposure. The paper reports component metrics so a reviewer can inspect the tradeoff rather than trusting the scalar alone.

\section{Generator Audit}

The deterministic generator maps each factor tuple to latent identity risk, geometry risk, re-acquisition need, safety pressure, visibility pressure, diagnosis quality, and active information. Each policy maps those latent quantities to component metrics. Small deterministic variation is derived from factor codes so reruns are reproducible.

This is not a physics simulator. It does not model contact dynamics, image formation, or object geometry. It is a mechanism audit: a way to stress the goal-identity decision independently from low-level robot dynamics.

The benefit of this design is inspectability. A reviewer can read the runner and see why ambiguity affects swaps, why horizon affects drift, why observability affects re-acquisition, and why cost affects search utility.

\section{Represented Scale}

The compact CSV stores 518,400 rows. Each row represents 19 seeds, 7 robot/object instances, 5 layouts, 4 disappearance schedules, 4 perception calibrations, and 32 episodes. Each episode has 80 planning ticks. This yields 340,480 represented evaluations per row and 27,238,400 represented planning-tick decisions per row.

The total represented evaluation count is 176,504,832,000. The total represented planning-tick decision count is 14,120,386,560,000. These are represented counts, not explicit physics ticks. The runner is RAM-light because it streams compact rows and keeps only summary accumulators.

\section{Reproducibility Map}

The full-scale runner is \path{scripts/run_full_scale_disappearing_goal_suite.py}. The row-level table is \path{results/full_scale/condition_metrics.csv}. Summary CSVs, LaTeX tables, validation JSON, and figures live in \path{results/full_scale/} and \path{figures/full_scale/}. The v2 negative control remains in \path{docs/v2_rebinding_ambiguity_stress.json} and \path{paper/v2_rebinding_ambiguity_table.tex}.

To reproduce the generated evidence:
\begin{verbatim}
python scripts/run_full_scale_disappearing_goal_suite.py
\end{verbatim}
Then build the manuscript from the paper directory or use the final build script.

\section{Real-Robot Validation Protocol}

A real validation should instrument a manipulation stack with goal-proxy logs. The logs should record the goal identity, geometry anchor, visibility status, re-binding gate, accepted observations, rejected observations, re-acquisition actions, search cost, final action, and whether the final action used the correct referent.

The task set should include self-occluded shelf insertion, hidden drawer placement, human handover, bimanual assembly, bin picking with distractors, language referent manipulation, and failed-grasp recovery. Each task should have clear, close-distractor, appearance-similar, moving-goal, and false-observation variants.

The report should compare visible-only, last-seen belief, loose proxy, proxy re-acquisition, active search, risk-calibrated proxy, and a hindsight human-labeled oracle. It should report the same component metrics used here.

\section{Reviewer Attack Matrix}

\paragraph{Attack: This is just partial observability.}
Response: partial observability is the setting. The contribution is the goal-specific proxy interface and identity re-binding audit.

\paragraph{Attack: The benchmark is synthetic.}
Response: correct. It is a mechanism benchmark, not deployment evidence.

\paragraph{Attack: Active search is enough.}
Response: active search is a strong baseline, but it pays high search and delay costs. The final result shows calibrated search has better utility.

\paragraph{Attack: The oracle is much better.}
Response: correct. The oracle is an upper bound and shows room for better identity estimation.

\paragraph{Attack: Loose proxy is unfair.}
Response: loose proxy is included because it was the failure found by v2. It represents the common mistake of accepting a nearby reappearing object without an identity gate.

\section{Claim Guardrails}

Use identity-calibrated goal proxies rather than occlusion solved. Use best non-oracle in a deterministic benchmark rather than state of the art. State that re-acquisition is necessary but not free. State that active search is strong but costly. State that real robot validation remains future work.

\section{Falsification Criteria}

The claim would be weakened if visible-only, last-seen belief, loose proxy acceptance, or active search matched \method across ambiguity, horizon, occlusion, observability, cost, and task regimes while keeping swaps and search cost low. It would also be weakened if real robots cannot log identity swaps reliably.

The claim would be strengthened by real robot experiments showing the same pattern: disappearing goals require persistent identity, loose re-binding causes swaps under distractors, active search reduces swaps but costs time, and calibrated re-acquisition gives the best tradeoff.

\section{Artifact Integrity Checklist}

The final artifact should pass six checks. The PDF should be at least 25 pages. The final PDF should be exported to \path{C:/Users/wangz/Downloads/56.pdf} only after finalization. Local \path{paper/main.pdf} should be removed. The generated validation JSON should match the represented counts in the paper. Current status docs should mark v3 final status, with v2 preserved only as a negative control. The final hash, file size, page count, and visual QA status should be recorded.

\section{Implementation Pattern: Task-And-Motion Planning}

Task-and-motion planners often separate symbolic task goals from geometric action planning. A disappearing-goal proxy can sit between those layers. The symbolic planner says that the goal is, for example, the left cup, the drawer interior, or the shelf slot. The geometric planner receives a proxy with the current anchor, uncertainty, affordance, and gate. When visibility disappears, the symbolic goal remains stable while the geometric anchor becomes uncertain.

The important engineering detail is that re-binding should not occur implicitly inside perception. If the perception module reports a candidate cup after occlusion, the planner should ask whether the candidate is compatible with the original proxy. Compatibility can use geometry, appearance, language context, contact history, and task role. If compatibility is weak, the policy should search or request confirmation before acting.

This interface makes failures legible. A task-and-motion planner can report that planning failed because the proxy was lost, because the re-binding gate rejected all observations, because search cost was too high, or because the action was safety critical. These are different failures and should not be collapsed into one generic goal-not-found state.

\section{Implementation Pattern: Behavior Trees}

Behavior trees already contain fallback nodes, retry decorators, condition checks, and recovery branches. A goal proxy can be stored on the tree blackboard. Action nodes read the proxy, condition nodes test its validity, and recovery nodes decide whether to search, reset, or escalate.

The usual retry decorator should be modified. It should not simply retry the same action while the target is hidden. After each failed action or disappearance event, it should query the proxy policy. If identity risk is low, the decorator may continue. If ambiguity is rising, it should run a re-acquisition subtree. If the gate rejects the reappearing observation, it should mark a proxy-loss failure rather than silently switching goals.

This is a low-friction integration path. The paper does not require replacing behavior-tree recovery. It requires exposing the goal identity and re-binding decision as first-class blackboard state.

\section{Implementation Pattern: Learned Policies}

Learned manipulation policies often encode goal information as image features, language embeddings, object tokens, or latent memory. A learned policy can still implement the proxy contract if the representation exposes identity confidence, stale age, accepted observations, rejected observations, and re-acquisition actions.

The final paper's benchmark is deterministic, but the logging contract applies to learned systems. If a transformer policy attends to an object token after occlusion, the system should record whether that token is still believed to refer to the original goal. If an object-centric model updates a slot, it should record whether the update passed an identity gate. If a language-conditioned model changes attention from one object to another, the report should count that as a possible swap.

This is important because learned policies can hide goal switches inside representation dynamics. A benchmark that reports only task success may reward a policy that completes the wrong task.

\section{Trace Example: Shelf Slot Insertion}

A shelf insertion task begins with a visible slot. As the robot approaches, the wrist and object occlude the slot. A visible-only controller loses the target. A last-seen controller continues toward a remembered location, but the final pose may be stale if the shelf or wrist alignment changed. A loose proxy controller may accept a nearby shelf opening as the old slot if it reappears near the remembered anchor.

\method preserves the slot identity while tracking geometry uncertainty. If wrist camera evidence is cheap, it performs a short active view motion. If the slot remains ambiguous, it tightens the gate and requests geometric confirmation. If the action is nearly irreversible, it stops rather than forcing insertion into a possibly wrong slot.

The key point is that the robot does not merely remember the slot. It knows whether the remembered proxy is still safe to use.

\section{Trace Example: Human Handover}

In a handover, the receiving pose can disappear behind the object or the robot's own hand. The identity of the goal is the human's receiving configuration, not just the last visible point. If the human shifts while hidden, last-seen belief can become dangerous. If another visible hand-like region appears, loose re-binding can be wrong.

\method raises identity risk when human-proximity and hidden motion are high. It may pause, ask for visual re-acquisition, wait for tactile or force confirmation, or abort the handover. This behavior can reduce task speed, but it prevents the worse failure: acting confidently on a stale human pose.

This trace explains why human-proximity regimes should carry a safety exposure metric. A policy with high success but high unsafe search or action exposure should not be treated as a good disappearing-goal policy.

\section{Trace Example: Language Referent Clearing}

Consider a command such as move the red mug to the tray. The red mug is visible at instruction time, then becomes hidden behind another object. A semantic anchor helps: the goal is not any cup, but the red mug that was selected by the instruction. However, semantic anchors are not perfect. If two red mugs are nearby, appearance alone may not resolve identity.

The proxy state should store the language role, object identity evidence, geometry, and ambiguity. If a red mug reappears, the gate should ask whether it matches the original mug, not whether it satisfies the language label in isolation. This distinction is often missing in high-level language-conditioned manipulation reports.

The benchmark's semantic-anchor observability regime is therefore not a claim that language solves disappearance. It is a stress axis showing that semantic evidence can help identity, but only when the gate treats it as evidence rather than proof.

\section{Trace Example: Bin Picking With Distractors}

Bin picking is a natural identity-swap trap. The target can be visible at the start, hidden during grasp approach, and replaced in the camera by a similar distractor. A loose proxy policy may accept whichever candidate appears closest to the old anchor. That can give low final geometric error relative to the distractor while failing the task.

The correct report must therefore separate final error from identity swap. A policy can have low final error and high swap rate if it accurately moves to the wrong object. This is why the full-scale table includes both error and swap.

For bin picking, active search may be justified because distractor risk is high. But if the bin task is low cost and many objects are acceptable substitutes, the utility weights might change. The paper's claim is about tasks where the original referent matters.

\section{Trace Example: Drawer Interior Placement}

Drawer placement illustrates container-interior disappearance. The target region is visible before approach, then hidden by the drawer geometry and the object being placed. A camera may not recover visibility without retreating, so tactile or geometric confirmation can be more useful than visual search.

\method can choose confirmation based on the observability regime. In a tactile/geometric regime, it may accept contact constraints as identity evidence. In a passive vision regime, it may not. This is why the benchmark separates re-acquisition modalities rather than treating search as one scalar.

The trace also shows why action cost matters. Pulling out and re-viewing the drawer may be expensive but safe. Inserting without confirmation may be fast but irreversible.

\section{Utility Sensitivity}

The final utility penalizes identity swaps strongly. This is not arbitrary. If the robot acts on the wrong referent, it has changed the task. A geometric success on the wrong object should not be rewarded as a normal near miss.

If the swap penalty were reduced, loose proxy acceptance would improve substantially because its main failure would be discounted. That would correspond to a domain where any similar object is acceptable. Such a domain is not the disappearing-goal problem studied here. It is a different task with fungible goals.

If the search penalty were raised, active search would become weaker and visible-only might look less bad. This could be appropriate for very time-critical domains. But even there, a policy should still report swaps and proxy loss. The component metrics are included so a reader can reinterpret the results under different priorities.

\section{Why Active Search Is Not Enough}

Active search is attractive because it reduces ambiguity. In the benchmark, active-search proxy has low final error, strong re-acquisition success, and low swap rate. It nearly matches \method on swaps. The problem is cost. Its aggregate search effort is 0.756, delay is 0.483, and waste is 0.277.

\method searches less because it conditions search on identity risk and cost. It accepts that some easy regimes do not need aggressive information gathering. This is the core engineering tradeoff: too little search causes swaps, while too much search wastes time and can reduce task utility.

A real robot paper should therefore report active search as a serious baseline. It should not compare only to no-memory or last-seen baselines.

\section{Why Visible-Only Is Still Useful}

Visible-only control is weak under disappearance, but it is not useless as a baseline. It has low swap rate because it refuses to act on hidden or ambiguous goals. That makes it a conservative lower bound on identity risk.

The visible-only failure is proxy loss. Its aggregate proxy loss is high because it lacks a persistent task carrier. This distinction clarifies the problem. A disappearing-goal policy should improve success and progress without inheriting the swap failures of loose memory.

The final paper uses visible-only to expose the cost of refusing to maintain identity through occlusion.

\section{Why Last-Seen Belief Fails}

Last-seen belief carries geometry but not enough identity. It keeps moving toward a remembered estimate, which can be helpful during short disappearance. Under long disappearance, moving targets, hidden motion, or false reappearance, the estimate drifts. When a distractor appears near the old anchor, the policy has no gate that says whether the new observation is the original goal.

The full-scale benchmark gives last-seen belief aggregate swap rate 0.392 and utility -0.259. This does not mean all belief-state methods fail. It means that a belief without an explicit identity and re-binding contract is not enough for disappearing manipulation goals.

\section{Why Loose Proxy Acceptance Is Dangerous}

Loose proxy acceptance appears more sophisticated than last-seen belief because it has a proxy. But the proxy accepts compatible-looking observations too easily. The result is worse than drift: the planner can become confident in the wrong referent.

This is why the paper preserves the v2 negative control. Under close distractors, loose proxy acceptance failed catastrophically. In the full-scale benchmark, its aggregate swap rate is 0.521 and utility is -0.261. A proxy interface is only useful if it can reject tempting but wrong re-bindings.

\section{Oracle Gap}

The oracle proxy has utility 0.999, success 0.971, swap rate 0.001, and re-binding precision 0.985. It is intentionally unrealistic because it knows true identity and ambiguity labels. Its purpose is to show the ceiling.

The gap between oracle and \method is useful. It shows that better identity estimation, better re-acquisition, and better learned proxy updates could improve the result. The paper should not hide this gap. It should use the gap to motivate real robot logging and future learned methods.

\section{Condition-Row Audit}

The row-level CSV lets a reviewer inspect any exact factor tuple. For example, they can select human handover, mixed long horizon, human occlusion, adversarial false observation, multimodal delayed confirmation, human-proximity safety cost, and active search. The row will expose success, error, swap, proxy loss, precision, re-acquisition, search, delay, waste, safety, progress, utility, and represented weight.

This audit trail matters because aggregate tables can hide tradeoffs. If a reviewer thinks \method wins only by under-searching, they can inspect cost-specific rows. If they think active search should win in high ambiguity, they can filter ambiguity regimes. If they think loose proxy is unfair, they can inspect clear reappearance rows where loose proxy has fewer disadvantages.

\section{Failure-Token Logging Schema}

A real robot should log goal-proxy events with a compact schema. Each event should include task identifier, goal identity, visibility state, geometry anchor, uncertainty, stale age, re-binding candidates, accepted observation, rejected observations, re-acquisition action, search cost, action selected, and final identity label.

The log should distinguish three failures. Proxy loss means the robot no longer has a usable goal carrier. Identity swap means the robot acted on a different referent. Re-acquisition failure means search did not recover enough evidence. These failures need separate counts.

This schema is simple enough for behavior-tree, TAMP, and learned-policy stacks. It gives reviewers the data needed to judge disappearing-goal behavior.

\section{Minimal Real-Robot Report Template}

A real robot report should include at least seven tables. First, aggregate policy comparison. Second, ambiguity stress. Third, disappearance-horizon stress. Fourth, occlusion-cause stress. Fifth, re-acquisition modality stress. Sixth, dynamics and cost stress. Seventh, task-family summary.

The report should include traces. A trace should show when the goal disappears, how the proxy updates, which observations are rejected, when search occurs, and whether the final action uses the correct referent. A single final success number is not enough.

The report should also show failures. A strong paper should include examples where loose proxy swaps, active search wastes time, \method stops correctly, and \method fails due to ambiguous evidence.

\section{Benchmark Extension Plan}

The deterministic benchmark can be extended in three ways. A physics extension could connect factor tuples to simulated camera occlusion and contact. A learned-policy extension could replace deterministic policy equations with learned proxy updates. A real-trace extension could replay logged robot episodes through the same metric pipeline.

Any extension should keep the same baselines. Visible-only, last-seen belief, loose proxy, re-acquisition, active search, \method, and oracle each expose a different failure mode. Removing one of them makes the benchmark easier but less informative.

The extension should also keep row-level output. A black-box benchmark without auditable rows would weaken the paper's reporting discipline.

\section{Practical Deployment Defaults}

For a robot team that cannot implement \method immediately, the benchmark suggests conservative defaults. Do not silently accept nearby reappearing objects. Log identity swaps separately from final error. Use active re-acquisition when ambiguity is high. Treat human-proximity disappearance as safety critical. Preserve a stale-age field on every goal proxy.

For a team that can implement risk calibration, start with a simple score. Increase identity risk for close distractors, appearance-similar objects, hidden motion, false observations, and irreversible actions. Decrease risk when tactile, geometric, semantic, or multimodal evidence supports the original identity. Tune the search threshold on held-out tasks.

\section{Final Manuscript Contract}

The final manuscript should answer five questions. What failed in v2? Loose proxy acceptance swapped to distractors. What changed in v3? The claim became identity-calibrated proxies with explicit re-acquisition and cost. Which baselines are serious? Proxy re-acquisition and active search. Where does the method degrade? Adversarial false observations, moving hidden goals, and high-cost human-proximity regimes. What is not claimed? Real-robot deployment safety.

\section{Handoff Notes}

Future continuation should treat this as the Paper56 v3 full-scale track, not the old v2 workshop note. The generated validation JSON is authoritative for represented scale. If the runner or manuscript changes, rebuild the PDF, rerender representative pages, update the hash metadata, and rerun stale-status scans before committing.

\section{Task Family Definitions}

The 12 task families are designed to cover distinct ways that a goal can disappear. Tabletop cup placement is the simplest case: the goal is usually stable, but the hand or carried object can hide it during approach. Shelf slot insertion adds tight geometry and a goal feature that can disappear behind the wrist. Drawer interior placement makes the goal a region rather than a standalone object. Human handover adds safety risk and hidden motion.

Bimanual covered assembly stresses self-occlusion because one hand can cover the target feature for the other. Tool alignment with a hidden fastener stresses fine geometry: a proxy must preserve both identity and affordance. Cable aperture routing combines deformable-object uncertainty with a hidden aperture. Bin picking target identity stresses distractors and false reappearance. Mobile manipulation view loss adds robot-base motion and viewpoint change.

Language referent clearing tests semantic anchors. Regrasp hidden target frame tests whether a proxy survives object-relative frame changes. Failed grasp original target tests recovery: after a failed action, the robot must know whether it is still trying to recover the original target or has drifted to a nearby object.

\section{Horizon Regime Definitions}

Short disappearance tests whether a policy can bridge a brief visibility gap. Medium disappearance tests approach-time memory. Long decisive disappearance is harder because the goal is hidden at the moment the action commits. Delayed execution adds stale time between observation and action. Intermittent receding-horizon visibility tests repeated disappear/reappear cycles. Mixed long horizon tests several disappearance intervals in one task.

The horizon axis is not just duration. It changes when evidence is available. A short disappearance may be solved by last-seen memory. A long decisive disappearance may require active confirmation before committing. A receding-horizon regime can repeatedly refresh or corrupt the proxy. A delayed execution regime tests whether stale age is reported.

\section{Occlusion Regime Definitions}

Robot self-occlusion occurs when the manipulator blocks its own camera. Carried-object occlusion occurs when the object in the gripper blocks the target. Environmental occlusion occurs when shelf, drawer, bin, or fixture geometry hides the goal. Human occlusion occurs when a person temporarily blocks the referent. Container-interior disappearance occurs when the target region is inside a container. Late full reappearance occurs when the goal returns after the policy has already made several decisions.

These regimes matter because the right re-acquisition action differs. A wrist camera helps self-occlusion. Geometric or tactile confirmation helps container interiors. Human occlusion may require waiting rather than searching. Late reappearance needs a gate because the robot may already have drifted.

\section{Ambiguity Regime Definitions}

Clear reappearance is the easiest regime because the original goal comes back with little distractor pressure. Close spatial distractor means another candidate appears near the old anchor. Appearance-similar distractor means visual identity is weak. Crossing trajectories mean the hidden goal and distractor can exchange apparent positions. Goal moves while hidden means last-seen geometry becomes stale. Adversarial false observation means the observation stream offers a tempting but wrong update.

The ambiguity axis is the main reason persistent memory is not enough. If the goal disappears and reappears alone, a proxy can be simple. If the goal reappears near a distractor, the proxy must decide whether to accept the observation. If the observation is false, a loose proxy becomes actively harmful.

\section{Observability Regime Definitions}

Passive vision-only reappearance gives no extra information beyond what becomes visible. Active wrist camera adds viewpoint control. Tactile/geometric confirmation adds contact or constraint evidence. Semantic anchor adds language or object-role evidence. Multimodal delayed confirmation combines several sources but costs time.

The benchmark treats these as regimes rather than one scalar because each supports different robot designs. A camera-only tabletop robot, a tactile insertion robot, and a language-conditioned mobile manipulator should not be evaluated as though they have the same re-acquisition evidence.

\section{Dynamics And Cost Regime Definitions}

Static cheap search is the easiest cost regime. Moving target moderate cost adds hidden motion. Moving distractor ambiguity adds false re-binding risk. High action cost makes search and delay expensive. Human-proximity safety critical makes identity errors and irreversible actions more costly.

The cost axis prevents a naive reading of the result. Active search can reduce swaps, but if search delays the task or creates safety exposure, it may not be optimal. \method wins by treating search as a budgeted decision.

\section{Policy Failure Modes}

Visible-only fails by proxy loss. It refuses to act without current observations, so it keeps identity risk low but loses tasks. Last-seen belief fails by drift and swap. It has a target, but not enough identity evidence. Loose proxy fails by over-acceptance. It updates too easily and can become confidently wrong.

Identity-gated proxy fails by being too passive. It rejects suspicious observations but often lacks a good recovery action. Proxy re-acquisition fails by spending search and still accepting some ambiguous cases. Active search fails by over-searching. \method can fail if its risk estimate is wrong, especially when ambiguity is high but evidence appears deceptively clean. The oracle fails only through the residual noise built into the deterministic generator.

\section{Reading The Main Table}

The main table should be read left to right. Success alone favors proxy re-acquisition, which completes many tasks. Error alone favors active search, which often gets better geometry. Swap rate penalizes loose proxy and last-seen belief. Re-acquisition success rewards active search and oracle. Precision rewards oracle and \method. Utility combines these under the paper's safety-sensitive priorities.

The key comparison is \method versus proxy re-acquisition and active search. Proxy re-acquisition has higher success but higher swaps and search. Active search has similar swaps but much higher delay and waste. \method gives the best utility because it balances these components.

\section{Reading The Ambiguity Table}

The ambiguity table should be read as a stress test, not just a performance table. The loose proxy column is the negative-control alarm. If that column did not rise with ambiguity, the benchmark would not be testing the right failure. The re-acquisition column shows that active recovery helps but does not eliminate swaps. The risk-calibrated column shows the desired behavior: lower swaps and useful success under all regimes, with degradation under adversarial false observations.

The strongest regime is adversarial false observation. It is not meant to be common. It asks whether the policy can avoid a tempting wrong update. A disappearing-goal method that fails there should say so.

\section{Reading The Horizon Table}

The horizon table tests whether the policy can preserve identity over time. Visible-only should degrade as disappearance grows. Last-seen error and swaps should rise. \method should degrade less sharply because it uses uncertainty, stale age, and re-acquisition.

If a future method has high performance under short disappearance but collapses under mixed long horizons, it has not solved the core problem. It has solved momentary occlusion.

\section{Reading The Cost Table}

The cost table compares active search and calibrated search. Active search spends high search effort in every regime. That is appropriate when ambiguity is severe and search is cheap, but wasteful when identity risk is low or search is expensive. \method searches less and keeps utility higher.

The human-proximity regime is important. A policy should be conservative near people, but it should not wander through search motions that increase exposure. Risk calibration must consider both identity and action safety.

\section{How To Instrument Proxy Loss}

Proxy loss should be logged when the system no longer has a usable goal carrier. This can happen because the goal disappeared too long, all re-binding candidates were rejected, geometry uncertainty exceeded a threshold, or the task role became ambiguous. Proxy loss is not the same as identity swap. In proxy loss, the robot knows it does not know. In identity swap, the robot acts on the wrong referent.

This distinction makes visible-only and loose proxy comparable. Visible-only has high proxy loss and low swaps. Loose proxy has low proxy loss and high swaps. A good policy should avoid both.

\section{How To Instrument Identity Swaps}

Identity swaps require ground-truth or human-labeled referent identity. In simulation, the label is known. On a real robot, the label may require manual annotation, fiducials, object instance tracking, or post-hoc video review. The report should state how swaps are labeled.

Swap metrics should be counted at the moment of action, not only at reappearance. A proxy may accept a wrong object and later recover before acting. That is a re-binding error but not necessarily a task swap. Conversely, a policy may never explicitly re-bind but still act on the wrong referent. The action-time label matters.

\section{How To Instrument Re-Binding Precision}

Re-binding precision measures whether accepted observations are correct. It should not be conflated with re-acquisition success. A policy can search actively and find candidates but accept too many wrong ones. Precision asks whether accepted updates preserve identity.

In a real system, precision can be estimated from annotated reappearance events. Each accepted candidate should be labeled as original goal, distractor, or ambiguous. Ambiguous labels should be reported separately rather than forced into success or failure.

\section{How To Instrument Search Effort}

Search effort includes camera motion, wrist motion, tactile probing, semantic questioning, additional perception calls, and waiting for a better view. The benchmark reports it as a scalar, but a real robot should break it into categories.

This matters because search effort has different costs. A brief wrist camera motion may be cheap. A tactile probe may be safe in insertion but unsafe in handover. Asking a human for clarification may be acceptable in collaboration but not in autonomous warehouse sorting.

\section{How To Instrument Safety Exposure}

Safety exposure counts risky action while identity is uncertain. It can include motion near a human, forceful insertion into a hidden slot, reaching into clutter with uncertain target identity, or moving an object toward an unconfirmed goal.

The paper does not make a formal safety proof. It asks that disappearing-goal studies report this exposure rather than hiding it behind success. In human-proximity tasks, identity uncertainty and physical risk interact.

\section{Relation To Object Permanence}

Object permanence says that objects continue to exist when hidden. Disappearing-goal manipulation needs something stricter: task permanence. The robot must preserve not only that an object exists, but that this object or region remains the intended goal of the current task.

An object permanence model can support the proxy, but it is not enough. If two objects are similar, both can be permanent. The planner needs to know which one is the task referent.

\section{Relation To Tracking}

Tracking estimates where an object goes. Goal proxies use tracking evidence but also carry task role and re-binding rules. A tracker may switch tracks under occlusion; a goal proxy should report that switch as an identity risk.

The full-scale benchmark is therefore not a tracking benchmark. It measures whether the manipulation planner's goal carrier remains valid for action.

\section{Relation To Active Perception}

Active perception gathers information. \method decides when information is worth the cost. The active-search baseline shows that more information gathering is not automatically better. A robot that always searches can be safe but inefficient; a robot that never searches can be fast but wrong.

The final claim is that disappearing-goal manipulation needs calibrated active perception tied to goal identity.

\section{Relation To Language Grounding}

Language can define a goal referent, but it can also be ambiguous. The red cup may identify one object at instruction time and several objects after occlusion. A language-conditioned policy should carry the original grounding as a proxy field, not re-ground the instruction from scratch after every disappearance.

This is especially important for long-horizon policies that use language repeatedly. Re-grounding without identity memory can silently change the task.

\section{Relation To World Models}

World models predict hidden state. A goal proxy can be represented inside a world model, but the planner still needs an action-facing interface. The question is not only whether the hidden goal state is predicted. The question is whether the predicted state still belongs to the intended referent.

The benchmark's last-seen and loose-proxy baselines are deliberately simple, but the metrics apply to learned world models as well.

\section{Why The Result Is Positive}

The final version has positive findings that the old short draft did not have. \method is the best non-oracle policy by utility. The oracle remains clearly better, preserving a meaningful upper bound. Proxy re-acquisition and active search are strong baselines, so the method is not winning against weak alternatives. Loose proxy acceptance and last-seen belief fail for different reasons, exposing the need for identity metrics.

The benchmark also makes progress measurable across several axes. Ambiguity raises swaps. Horizon length raises visible-only loss and last-seen drift. Cost regimes expose over-search. Observability regimes expose which evidence is available for re-binding. Task summaries show that the effect is not isolated to one scenario.

\section{What Would Make The Paper Stronger}

The paper would be stronger with real robot logs, object-instance annotations, and physical occlusion trajectories. It would be stronger with learned proxy updates trained on a subset of tasks and evaluated on held-out task families. It would be stronger with human-labeled ambiguity and safety exposure.

These are future directions, not hidden claims. The final manuscript states the limitation and provides a protocol for addressing it.

\section{Final Acceptance Checklist}

Before final export, verify the following. The runner must produce exactly 518,400 compact rows. The represented evaluation count must be 176,504,832,000. The represented planning-tick count must be 14,120,386,560,000. The proposed method must be best non-oracle, with the oracle best overall. The manuscript must be at least 25 pages. The final PDF must be exported to Downloads only after finalization. The local PDF must be removed. The final hash must be recorded.

\section{Submission Claim Matrix}

The final paper should make four positive claims. First, disappearing-goal manipulation is an identity problem, not just an occlusion problem. Second, persistent proxies help only when re-binding gates and re-acquisition actions are explicit. Third, in the deterministic full-scale benchmark, \method is the best non-oracle policy by utility. Fourth, the benchmark provides a reporting discipline for future real robot studies.

The final paper should avoid four stronger claims. It should not claim that \method solves partial observability. It should not claim deployment-ready safety. It should not claim that the risk score is the only possible policy. It should not claim that synthetic represented ticks are physical robot trials.

This matrix keeps the paper ambitious without becoming brittle. The method has a real positive result, but the scope is carefully bounded.

\section{Reviewer Response Catalog}

\paragraph{This is just memory.}
The response is that memory without identity gating fails. Last-seen belief has swap rate 0.392 and utility -0.259. Loose proxy memory has swap rate 0.521 and utility -0.261. The method's contribution is calibrated identity re-binding, not memory alone.

\paragraph{This is just active perception.}
Active perception is a strong baseline. Active-search proxy has low error and low swaps, but search effort 0.756 and utility 0.453. \method uses less search, 0.309, and has higher utility, 0.595. The difference is cost-aware identity calibration.

\paragraph{This is synthetic.}
Correct. The manuscript claims a deterministic mechanism benchmark and reporting discipline. It gives a real-robot validation protocol rather than pretending the synthetic benchmark is deployment evidence.

\paragraph{The oracle is much better.}
Correct. The oracle is a ceiling and sanity check. The gap shows room for learned identity estimation and better re-acquisition.

\paragraph{Success is not highest for the proposed method.}
Correct. Proxy re-acquisition has higher aggregate success, but also higher swaps, search, and delay. The paper's objective is not raw success at any identity cost. It is correct-goal manipulation under disappearance.

\section{Ablations To Add In A Real Follow-Up}

A real follow-up should ablate the proxy fields. Remove semantic anchors and measure language-referent swaps. Remove stale age and measure long-horizon drift. Remove rejected-observation logging and measure hidden swap events. Remove tactile/geometric confirmation and measure drawer or insertion errors. Remove action-cost awareness and measure over-search.

The follow-up should also ablate supervision. A learned proxy update trained with only final success may learn to swap when distractors are acceptable in training. Adding identity labels should reduce swaps. Adding rejected-candidate labels should improve precision. Adding action-cost labels should reduce wasted search.

These ablations are not required for the deterministic paper, but they are the path from mechanism benchmark to robot study.

\section{Dataset Card For The Generated Benchmark}

The generated benchmark is deterministic. It contains 518,400 compact rows. Each row is a factor tuple, not a physical rollout. Rows represent repeated seeds, instances, layouts, disappearance schedules, calibrations, episodes, and ticks. The data should be used to audit mechanism tradeoffs, not to train a robot policy.

The benchmark covers disappearing-goal regimes that are easy to miss in standard manipulation evaluations. It over-represents ambiguity and false observation on purpose because the goal is to stress identity maintenance. A real deployment distribution may contain fewer adversarial cases, but a submission-grade method should survive them.

Known limitations include synthetic formulas, no image formation, no contact dynamics, no real annotation noise, and no learned representation. Known strengths include full factor coverage, row-level auditability, deterministic reproducibility, serious baselines, oracle ceiling, and explicit negative controls.

\section{Paper-To-Artifact Traceability}

Every major claim in the paper maps to an artifact. The aggregate policy claim maps to \path{results/full_scale/policy_summary.csv} and \path{table_main_performance.tex}. The ambiguity claim maps to \path{ambiguity_policy_summary.csv} and \path{table_ambiguity_stress.tex}. The horizon claim maps to \path{horizon_policy_summary.csv}. The cost claim maps to \path{cost_policy_summary.csv}. The represented scale maps to \path{experiment_validation.json}. The final PDF metadata maps to the build-status JSON and final validation JSON.

This traceability is part of the submission hardening. A reviewer should be able to move from prose to table to CSV to runner without trusting hand-entered numbers.

\section{Final Readiness Rationale}

The final version is submission-ready as a benchmark paper because it has a corrected claim, a serious negative control, full-scale deterministic coverage, strong baselines, an oracle ceiling, component metrics, visualized stress axes, reproducible artifacts, and explicit limitations. The old short manuscript had an idea. The final manuscript has an auditable result.

The paper remains honest about its boundary. Real robot validation is future work. The value of the final artifact is that it makes that future work precise: log identity, proxy state, re-binding decisions, re-acquisition actions, swaps, proxy loss, and search cost.
